\section{Periods and the relation lattice}\label{sec:periods}

This section proves Theorem~\ref{thm:main} and Corollary~\ref{cor:smooth}. Theorems~\ref{thm:real} and~\ref{thm:lift}
give the inclusion $K\subseteq\Rel(\delta)$ and item~(2) directly (Section~\ref{ssec:relations}). The reverse inclusion
is a period computation, which follows the mechanism of Greene, Morrison and Strominger \cite[\S2]{GMS95} inside the
polynomial family. Near the degenerate locus the period $\wp_q$ of a holomorphic volume form over $\delta_q$ is the
critical value $\varepsilon_P$ at $q_P$ times a holomorphic function $\mathcal B_q$ that has no zero near the locus
$\{\varepsilon_P=0\}$ (Lemma~\ref{lem:period}). A
relation among the classes $[\delta_q]$, even a relation modulo torsion, therefore gives the same relation among the
functions $\wp_q$ on a whole neighbourhood $U$ of the member, and among their differentials at the point $c^{(0)}$ of
$\mathcal Z$ below it. There the matrix of differentials is $\operatorname{diag}(\mathcal B_q)\,\mathbf R$ up to an error
$O(t^{\theta_0/2})$, where $\mathbf R$ is the integer matrix with rows $r_P$, whose left kernel is $K\otimes\Q$. So the
relations have rank at most $\operatorname{rank}K$, and saturation of $K$ gives equality over $\Z$ and excludes torsion.
The same matrix, with a contraction argument and the holomorphic implicit function theorem, describes the nodal locus
(Section~\ref{ssec:nodal}). Section~\ref{ssec:lagrangian} replaces the spheres $\delta_q$ by pairwise disjoint Lagrangian
spheres, and Section~\ref{ssec:smooth} proves Corollary~\ref{cor:smooth}.

Throughout, the configuration of Section~\ref{ssec:member} is fixed: $\mathcal T$ and $\check h$, the pulling
refinement $\mathcal F_C$ and the height $h_C$, the vectors $w_u$, the integer $k_0$ and $H=k_0\varphi+h_C$. We use the
notation of Section~\ref{sec:regime}: the compact set $\mathscr C_0\subset\mathcal Z$, the radius $r\le r_1$ and the
bound $\eta_0\le r^2/20$ of Definition~\ref{def:window}, the constant $C_\eta$ of Lemma~\ref{lem:eta}, the sets
$U_\C\subseteq U^+_\C$ of Definition~\ref{def:window}, the max-norm $\|\xi\|=\max_n|\xi_n|$ on $\C^{\mathcal A}$ with
$\|r_P\|_*=4$, and $\eta_P(\xi)=r_P\cdot\xi$, extended complex-linearly. A coefficient vector near $c^{(0)}$ is written
$c^{(0)}e^\xi$ with $\xi\in\C^{\mathcal A}$. Constants called $C$ depend only on the configuration and on
$\mathscr C_0$; $C(\eta_0)$ depends in addition on $\eta_0$. All thresholds on $t$ below depend only on these data.

\subsection{The neighbourhood \texorpdfstring{$U$}{U} and the periods}\label{ssec:U}

Fix $t$ below the threshold of Lemmas~\ref{lem:W}, \ref{lem:NM} and~\ref{lem:TR}, and a member $c=c^{(0)}e^{\xi_c}$ as
in Theorem~\ref{thm:main}: $c^{(0)}\in\mathscr C_0$, and $\xi_c$ real and in the window. For a node $q$, with
parallelogram $P$, and a coefficient vector $c'$ at which Lemma~\ref{lem:NM} applies, $q_P(c')$ and $\varepsilon_P(c')$
denote the Morse point of $f^{(1)}$ in $\mathbb D_P$ and its critical value. Below, sets of nodes are indexed by the
nodes $q$ or by their parallelograms $P$ interchangeably.

\begin{lemma}[the neighbourhood $U$]\label{lem:U}
Put $\eta_0':=3\eta_0/2$. The open set
\[
U:=\bigl\{c^{(0)}e^{\xi}:\ \xi\in\C^{\mathcal A},\ \|\xi\|<3/2,\ |\eta_P(\xi)|<\eta_0'\ \text{for every }P\bigr\}
\]
of complex coefficient vectors has the following properties.
\begin{enumerate}
\item $U_\C\subseteq U\subseteq U^+_\C$; in particular $U$ contains $c$ and $c^{(0)}$. The set $U$ is simply connected.
\item The conclusions of Lemmas~\ref{lem:W}(2), \ref{lem:NM} and~\ref{lem:TR}(2) hold for every $c'\in U$.
\item Each $\varepsilon_P$ is holomorphic on $U$, and on all of $U$
\begin{equation}\label{eq:depsilon}
d\varepsilon_P=e^{-\eta_P}\,r_P\cdot d\xi+O\bigl(C(\eta_0)\,t^{\theta_0/2}\bigr).
\end{equation}
\item For $c'\in U$ the singular points of $Y_{t,c'}$ are the points $q_P(c')$ with $\varepsilon_P(c')=0$, and each of
them is an ordinary double point.
\end{enumerate}
\end{lemma}

\begin{proof}
The defining bounds of $U_\C$, $U$ and $U^+_\C$ are $(1,\eta_0)$, $(3/2,3\eta_0/2)$ and $(2,2\eta_0)$, so
$U_\C\subseteq U\subseteq U^+_\C$. The map $\xi\mapsto c^{(0)}e^\xi$ is injective on $\{\|\operatorname{Im}\xi\|<\pi\}$, and
$3/2<\pi$; so $U$ is the image of a convex set under a homeomorphism onto its image, hence contractible. This proves
(1). Lemmas~\ref{lem:W}(2), \ref{lem:NM} and~\ref{lem:TR}(2) are stated for every coefficient vector in $U^+_\C$, for $t$
below the threshold fixed above, hence they hold on $U$. This proves (2).

The critical point $q_P(c')$ is nondegenerate, so by the holomorphic implicit function theorem it depends
holomorphically on $c'$, and so does $\varepsilon_P(c')=f^{(1)}(q_P(c'))$. By Lemma~\ref{lem:NM}(2) on $U^+_\C$,
$\varepsilon_P+\kappa_P=O(t^{\theta_0/2})$ on $U^+_\C$,
where $\kappa_P=e^{-\eta_P}-1$. Every point of $U$ is the centre of a polydisc in $\xi$ of polyradius
$\varrho_*:=\min(1/2,\eta_0/8)=\min\bigl(1/2,\eta_0/(2\|r_P\|_*)\bigr)$ contained in $U^+_\C$: on it
$\|\xi\|<3/2+1/2$ and $|\eta_P|<3\eta_0/2+\|r_P\|_*\varrho_*\le2\eta_0$. The Cauchy estimates on these polydiscs bound
$d(\varepsilon_P+\kappa_P)$ by $C(\eta_0)t^{\theta_0/2}$ on $U$, and $d\kappa_P=-e^{-\eta_P}\,r_P\cdot d\xi$. This is
(3). By Lemma~\ref{lem:W}(2) on $U$ the singular points of $Y_{t,c'}$ are the $q_P(c')$ with $\varepsilon_P(c')=0$. At
such a point the local defining function $f^{(1)}$ in the ambient chart has a nondegenerate critical point with critical
value $0$, so the point is an ordinary double point of $Y_{t,c'}$. This is (4).
\end{proof}

For a node $q$ with parallelogram $P$ put $D_q:=\{c'\in U:\varepsilon_P(c')=0\}$, and for $Q'\subseteq Q$ put
$D_{Q'}:=\bigcap_{q\in Q'}D_q$. By Lemma~\ref{lem:U}(4), $D_{Q'}$ is the locus of members of $U$ with a node at
$q_P$ for every $q\in Q'$. By~\eqref{eq:depsilon}, and since $|e^{-\eta_P}|\ge e^{-\eta_0'}\ge e^{-1}$ on $U$ while the
linear form $r_P\cdot d\xi$ has norm $\|r_P\|_*=4$, the differential $d\varepsilon_P$ does not vanish on $U$ for $t$ small;
so $D_q$ is a smooth complex hypersurface of $U$, possibly empty. The set $U^\circ:=U\setminus\bigcup_qD_q$ is connected,
as the complement of a proper analytic subset of a connected complex manifold, and it contains $c$, because
$\varepsilon_P(c)\neq0$ in the window (Lemma~\ref{lem:NM}(3)). The members of $U^\circ$ are smooth.

\subsubsection*{Vanishing spheres over $U$}
For $c'\in U\setminus D_q$ let $\delta_q(c')\subset Y_{t,c'}$ be the vanishing sphere of the Morse point $q_P(c')$ for the
straight path of values from $\varepsilon_P(c')$ to $0$, as in Definition~\ref{def:deltaq}: in holomorphic coordinates
$\zeta$ centred at $q_P(c')$ with $f^{(1)}=\varepsilon_P+\sum_i\zeta_i^2$, it is
$\{\zeta\in(-\varepsilon_P)^{1/2}\R^4:\sum_i\zeta_i^2=-\varepsilon_P\}$. On $U$,
\begin{equation}\label{eq:epsbound}
|\varepsilon_P|\le|\kappa_P|+Ct^{\theta_0/2}\le e^{\eta_0'}\eta_0'+Ct^{\theta_0/2}<r^2/6,
\end{equation}
since $\eta_0'=3\eta_0/2\le3r^2/40$; so the sphere lies in $\mathbb D_P$. We orient $\delta_q(c')$ by continuous
transport from the member $c$, at which it carries the orientation of Definition~\ref{def:orient}. This is well defined
on all of $U\setminus D_q$: since $U$ is simply connected and $D_q$ is a smooth hypersurface, $\pi_1(U\setminus D_q)$ is
generated by conjugates of meridians of $D_q$, and along a meridian $\varepsilon_P$ turns once around $0$, which carries
$\delta_q$ to itself with its orientation (Remark~\ref{rem:kappa}). The arguments below use Definition~\ref{def:orient}
only at $c$.

\begin{lemma}[flatness]\label{lem:flat}
The members $Y_{t,c'}$, $c'\in U^\circ$, form a locally trivial fibre bundle over $U^\circ$, and the classes
$[\delta_q(c')]\in H_3(Y_{t,c'};\Z)$, $q\in Q$, are flat sections of the local system of the groups
$H_3(Y_{t,c'};\Z)$ over $U^\circ$.
\end{lemma}

\begin{proof}
The set $\mathcal Y^\circ:=\{(c',y):c'\in U^\circ,\ y\in Y_{t,c'}\}$ is a closed submanifold of $U^\circ\times X_{\mathcal T}$,
and the projection to $U^\circ$ is a proper submersion, since every fibre is a smooth hypersurface of the compact
manifold $X_{\mathcal T}$. By Ehresmann's theorem it is locally trivial, and parallel transport along a path in $U^\circ$
is induced by a diffeomorphism of fibres, well defined up to isotopy. Over a small ball in $U^\circ$ the Morse points,
the Morse coordinates of the holomorphic Morse lemma with parameters \cite[\S\S6.2, 9.6]{AGZV1} and the spheres $\delta_q(c')$ vary
continuously, with the continuous orientation above; in a trivialisation over the ball the spheres form an isotopy, so
their classes correspond under parallel transport.
\end{proof}

\subsubsection*{Periods}
Let $\Theta_T:=\chi^{-n_1}d\chi^{n_1}\wedge\dots\wedge\chi^{-n_4}d\chi^{n_4}$ be the invariant four-form of the dense
torus for a basis $n_1,\dots,n_4$ of $N$, and write $F_t/x^0=\sum_n\gamma_n\chi^n$. In the chart of a maximal cone, with
Cox coordinates $Z_1,\dots,Z_4$, we have $\Theta_T=\pm dZ_1\wedge\dots\wedge dZ_4/(Z_1\cdots Z_4)$ and
$F_t/x^0=F_\sigma/(Z_1\cdots Z_4)$, where $F_\sigma$ is the defining polynomial of $Y_{t,c'}$ in the chart. So
$\Theta_T/(F_t/x^0)=\pm dZ_1\wedge\dots\wedge dZ_4/F_\sigma$ has a simple pole along $Y_{t,c'}$ and no other pole. For
$c'\in U$ its residue $\Omega_{t,c'}$ is a nowhere vanishing holomorphic three-form on the smooth locus of $Y_{t,c'}$,
depending holomorphically on $c'$. Put
\[
\wp_q(c'):=\int_{\delta_q(c')}\Omega_{t,c'},\qquad c'\in U\setminus D_q .
\]

\begin{lemma}[periods]\label{lem:period}
Let $q$ be a node, with parallelogram $P$. There are a holomorphic function $\mathcal B_q$ on $U$, a sign $s_q$ and a
constant $C$ independent of $t$ such that $\wp_q=\mathcal B_q\varepsilon_P$ on $U\setminus D_q$ and
\[
\bigl|\mathcal B_q-4\pi^2s_q\,\mu_P/c_0\bigr|\le C\bigl(|\varepsilon_P|+t^{\theta_0/2}\bigr)\qquad\text{on }U.
\]
In particular $\wp_q$ extends holomorphically to $U$, $\mathcal B_q$ is bounded on $U$ uniformly in $t$, and there is
$\varepsilon_*>0$ independent of $t$ such that, for $t$ small, $\mathcal B_q$ has no zero at the points of $U$ with
$|\varepsilon_P|\le\varepsilon_*$, which include $D_q$.
\end{lemma}

\begin{proof}
\emph{Coordinates.} Fix $c_*\in U$ and freeze the positive scale $\nu_*:=\nu_Z(c_*)$. On a neighbourhood
$V\subset U$ of $c_*$ use $\mu_P^{1/2}=\exp(\eta_P/2)$ and put $z_1:=Z/\nu_*$ and $z_2:=\mathfrak b(c')\nu_*Z'$. These coordinates depend holomorphically
on the point and on $c'$, and $z_1z_2=\mathfrak bZZ'$. At $c_*$ they agree with $\hat Z$ and
$\vartheta_P\hat Z'$, since $\nu_Z\nu_{Z'}=1/|\mathfrak b|$. The scales defining $\mathbb D_P$ are continuous in
$c'$, so, after shrinking $V$, these charts have a common domain inside $2\mathbb D_P$.
By Lemma~\ref{lem:NM}(1), at $c_*$ the function $f^{(1)}$ differs by $O(t^{\theta_0/2})$ from
$f_{\rm lin}:=xy+z_1z_2-\kappa_P$, since $a\hat Z$ and $a'\hat Z'$ are
$O(\lambda_P^{-1/2})$ there. Fibre-variable Cauchy estimates give the same bound in $C^3$ on $\mathbb D_P$;
on a slightly smaller common domain it holds throughout $V$. Uniform bounds below are taken at the arbitrary
centre $c_*$, where the original normalisation applies, so their constants are independent of $c_*$ and $t$. Let
$\zeta^0$ be the linear coordinates with $x=\zeta^0_1+i\zeta^0_2$, $y=\zeta^0_1-i\zeta^0_2$, $z_1=\zeta^0_3+i\zeta^0_4$,
$z_2=\zeta^0_3-i\zeta^0_4$, so that $f_{\rm lin}=\sum_i(\zeta^0_i)^2-\kappa_P$ and
$dx\wedge dy\wedge dz_1\wedge dz_2=-4\,d^4\zeta^0$, where $d^4\zeta:=d\zeta_1\wedge\dots\wedge d\zeta_4$ for any coordinates
$\zeta$. Put $\zeta^1:=\zeta^0-\zeta^0(q_P)$ and $\mathsf H(\zeta^1):=\int_0^1(1-s)\operatorname{Hess}f^{(1)}(s\zeta^1)\,ds$, the
Hessian taken in the coordinates $\zeta^1$. Since $q_P$ is a critical point of $f^{(1)}$ with value $\varepsilon_P$,
Taylor's formula with integral remainder gives $f^{(1)}=\varepsilon_P+(\zeta^1)^{\mathsf T}\mathsf H(\zeta^1)\,\zeta^1$. The Hessian
of $f_{\rm lin}$ is $2I$, and $f^{(1)}-f_{\rm lin}$ is $O(t^{\theta_0/2})$ in $C^3$ on the ball $\{|\zeta^1|<r/2\}$, which
lies in $\mathbb D_P$ because $r\le1/2$; so $\mathsf H=I+O(t^{\theta_0/2})$ in $C^1$ there. Put $\zeta:=\mathsf H(\zeta^1)^{1/2}\zeta^1$,
with the principal square root of the symmetric matrix $\mathsf H$. The critical point is holomorphic in $c'$ in
the original toric coordinates, by the implicit function theorem. Thus $\zeta$ are jointly holomorphic coordinates
on the local parameter chart, centred at $q_P$, with $f^{(1)}=\varepsilon_P+\mathfrak q$, $\mathfrak q:=\sum_i\zeta_i^2$, and
$\zeta-\zeta^1=O(t^{\theta_0/2})$ in $C^1$ on the ball, with uniform constants at $c_*$.
After shrinking $V$ and decreasing $t$, the image contains $|\zeta|<r/\sqrt5$.

\emph{The form.} On the chart $U_{u,u'}$ of Section~\ref{ssec:morse} we take $n_1,\dots,n_4=b-a,\,d-a,\,w_u,\,w_{u'}$.
Since $x^a/x^0=\chi^a=(ZZ')^{-1}$ by~\eqref{eq:nodeexp}, $\Theta_T/(F_t/x^0)=\Theta/f^{(1)}$ with
$\Theta:=\pm\gamma_a^{-1}(\chi_1\chi_2)^{-1}d\chi_1\wedge d\chi_2\wedge dZ\wedge dZ'$. From
$\chi_1=(\gamma_a/\gamma_b)\mu_P^{-1/2}(x-\mu_P^{-1/2})$, $\chi_2=(\gamma_a/\gamma_d)\mu_P^{-1/2}(y-\mu_P^{-1/2})$,
$Z=\nu_*z_1$, $Z'=z_2/(\mathfrak b\nu_*)$ and $\mathfrak b=\gamma_0/\gamma_a$, with
$\gamma_0=c_0$ because $H(0)=0$,
\[
\Theta=\pm\frac{\mu_P}{c_0}\,\frac{dx\wedge dy\wedge dz_1\wedge dz_2}{(1-\mu_P^{1/2}x)(1-\mu_P^{1/2}y)}
=\mp\frac{4\mu_P}{c_0}\,\frac{d^4\zeta^0}{(1-\mu_P^{1/2}x)(1-\mu_P^{1/2}y)} .
\]
Write $\Theta=g\,d^4\zeta$. Then $g$ is holomorphic in $c'$ and $\zeta$, bounded on the ball uniformly in $t$, and
$g(0)=\mp4(\mu_P/c_0)\bigl(1+O(t^{\theta_0/2})\bigr)$, because $q_P$ lies within $O(t^{\theta_0/2})$ of the node point
$x=y=0$ (Lemma~\ref{lem:NM}(2)).

\emph{The period.} For $|\lambda|<r^2/5$ let $\mathbf B_\lambda:=\lambda^{1/2}\mathbf B$, where $\mathbf B$ is the closed
unit ball of $\R^4\subset\C^4$ in the coordinates $\zeta$, oriented by $\R^4$; this does not depend on the choice of the
square root, since $-1$ preserves the orientation of $\R^4$. Its boundary lies in $\{\mathfrak q=\lambda\}$, and for
$\lambda=-\varepsilon_P$ it is the sphere $\delta_q(c')$; by~\eqref{eq:epsbound} this applies throughout $V$. Let
$s'_q=\pm1$ compare the orientation of $\delta_q$ with the boundary orientation of $\mathbf B_{-\varepsilon_P}$; it is
locally constant in $c'$. Put
\[
\mathcal I(c',\lambda):=\int_{\mathbf B_\lambda}g\,d^4\zeta=\lambda^2\int_{\mathbf B}g(\lambda^{1/2}u)\,d^4u .
\]
In the Taylor expansion of $g$ at $\zeta=0$ the terms of odd degree integrate to zero over $\mathbf B$, so $\mathcal I$
is holomorphic in $(c',\lambda)$, and $\mathcal I=\lambda^2\bigl(\tfrac{\pi^2}2g(0)+\lambda\mathcal I_1\bigr)$, since
$\mathbf B$ has volume $\pi^2/2$; by the Cauchy estimates $\mathcal I_1$ is bounded on $|\lambda|\le r^2/6$ uniformly in
$t$. On $Y_{t,c'}$ near $q_P$ the residue of $\Theta/f^{(1)}$ is $g$ times the Gelfand--Leray form $d^4\zeta/d\mathfrak q$,
and the Gelfand--Leray formula \cite{AGZV2} gives
$\partial_\lambda\mathcal I(c',\lambda)=\int_{\partial\mathbf B_\lambda}g\,d^4\zeta/d\mathfrak q$. Hence
\[
\wp_q=s'_q\,\partial_\lambda\mathcal I(c',-\varepsilon_P)
=-s'_q\,\varepsilon_P\bigl(\pi^2g(0)-\varepsilon_P\mathcal I_2(c',-\varepsilon_P)\bigr),\qquad
\mathcal I_2:=3\mathcal I_1+\lambda\,\partial_\lambda\mathcal I_1 .
\]
The function $\mathcal B_q:=-s'_q\bigl(\pi^2g(0)-\varepsilon_P\mathcal I_2(c',-\varepsilon_P)\bigr)$ is holomorphic near every
point of $U$, also on $D_q$, and it equals $\wp_q/\varepsilon_P$ on the dense set $U\setminus D_q$; so these local
functions define a holomorphic function on $U$. The estimate follows from the formula for $g(0)$, with
$s_q=\pm s'_q$. The sign agrees at each centre with the one in the continuously varying original normalisation,
so it is locally constant on the connected set $U\setminus D_q$ and hence constant.
Since $|\mu_P/c_0|$ is bounded below on $U$ by a constant depending only on $\mathscr C_0$ and $\eta_0$, the last
statement follows.
\end{proof}

\subsection{Relations: proof of Theorem \texorpdfstring{\ref{thm:main}}{A}(1) and (2)}\label{ssec:relations}

\subsubsection*{The threshold}
We use the choices of Section~\ref{ssec:member}: $r=r_1$, $\eta_1=r_1^2/20$ and $C=C_\eta$ in the statement of
Theorem~\ref{thm:main}, so that for $\eta_0\in(0,\eta_1]$ the members of Theorem~\ref{thm:main} are exactly the members
$c=c^{(0)}e^{\xi_c}$ with $c^{(0)}\in\mathscr C_0$ and $\xi_c$ real and in the window; here $\log\mu_P(c)=\eta_P(\xi_c)$.
Let $t_0$ be the threshold of Theorem~\ref{thm:lift} (Proposition~\ref{prop:Bprecise}) for $r$, $\eta_0$ and
$\mathscr C_0$, which lies below those of Lemmas~\ref{lem:W}, \ref{lem:NM} and~\ref{lem:TR}. In the proofs of items~(1)
and~(3) we decrease $t_0$ finitely many times, each time by a condition that involves only the configuration,
$\mathscr C_0$ and $\eta_0$; the result is the $t_0$ of Theorem~\ref{thm:main}. Fix $t<t_0$, a member $c$ and the
neighbourhood $U$ of Lemma~\ref{lem:U}.

\subsubsection*{Proof of item (2), and $K\subseteq\Rel(\delta)$}
Let $\mathfrak K_0$ be the barycentric subdivision of $\mathscr P$, in which $\Gamma$ is a full subcomplex
(Lemma~\ref{lem:full}), and let $k\in K$. By Theorem~\ref{thm:real} there is $z\in Z(\mathfrak K_0)$ with $\rho(z)=k$. By
Theorem~\ref{thm:lift}, applied to $\mathfrak K_0$ and $z$, there is an integral four-chain $\check\Xi(z)$ on $Y_t$ with
$\partial\check\Xi(z)=\sum_q\rho_q(z)[\delta_q]=\sum_qk_q[\delta_q]$. This is item~(2), and it gives
$\sum_qk_q[\delta_q]=0$ in $H_3(Y_t;\Z)$, so $K\subseteq\Rel(\delta)$.

Both theorems are stated with the same objects: the group $Z(\mathfrak K)$ of Definition~\ref{def:relcycle}, with
coefficients in the sheaf $\iota_*\Lambda$ and boundary in $C_1(\Gamma;\mathcal D)$, and the node coefficients of
Definition~\ref{def:nodecoef}. Theorem~\ref{thm:lift} reads $\rho_q$ through the passage formula $\rho_q=c_{ab}-c_{bc}$,
which agrees with Definition~\ref{def:nodecoef} (Section~\ref{ssec:cycles-bdry}), and it orients $\delta_q$ by
Definition~\ref{def:orient}, as Theorem~\ref{thm:main} does, for the same labelling of each parallelogram. So no further
compatibility is needed; by Lemma~\ref{lem:diag} the conclusion does not depend on the labelling.

\subsubsection*{Proof of item (1)}
Let $K_{\rm tor}:=\{k\in\Z^Q:\sum_qk_q[\delta_q]\in\Tors H_3(Y_t;\Z)\}$. Then $K\subseteq\Rel(\delta)\subseteq K_{\rm tor}$, and item~(1) is the
statement $K_{\rm tor}=K$. Let $\mathbf R$ be the $|Q|\times|\mathcal A|$ integer matrix with rows $r_P=e_a+e_c-e_b-e_d$; its entries lie in
$\{0,\pm1\}$. The points of the set $R$ of Section~\ref{ssec:admissible} are lattice points of $\partial\Delta$, so
$\Z^R\subseteq\Z^{\mathcal A}$, and under this inclusion $r_P=-\circv_q$, both for the labelling that orients $\delta_q$
(a change of diagonal reverses $r_P$, $\circv_q$ and $\delta_q$ together, Lemma~\ref{lem:diag}). Hence the left kernel of $\mathbf R$ over $\Q$ is
$\{k\in\Q^Q:\sum_qk_q\circv_q=0\}=K\otimes\Q$, and
\[
r_{\mathbf R}:=\operatorname{rank}\mathbf R=|Q|-\operatorname{rank}K .
\]

\emph{Step 1: relations propagate over $U$.} Let $k\in K_{\rm tor}$. By Lemma~\ref{lem:flat}, parallel transport along a path in
the connected set $U^\circ$ from $c$ to $c'$ is a group isomorphism $H_3(Y_{t,c};\Z)\to H_3(Y_{t,c'};\Z)$ carrying
$[\delta_q(c)]$ to $[\delta_q(c')]$; so $\sum_qk_q[\delta_q(c')]$ is a torsion class for every $c'\in U^\circ$.
Integration of the closed form $\Omega_{t,c'}$ is a homomorphism $H_3(Y_{t,c'};\Z)\to\C$, which vanishes on torsion
classes; hence $\sum_qk_q\wp_q=0$ on $U^\circ$. By Lemma~\ref{lem:period} each $\wp_q=\mathcal B_q\varepsilon_P$ extends
holomorphically to $U$, and $U^\circ$ is dense in $U$; so
\begin{equation}\label{eq:periodrel}
\sum_qk_q\,\wp_q\equiv0\quad\text{on }U,\qquad\text{for every }k\in K_{\rm tor},\ \text{in particular for every }k\in K .
\end{equation}

\emph{Step 2: the differentials at $c^{(0)}$.} At the point $c^{(0)}\in U$, where $\xi=0$, every $\eta_P$ vanishes, so
$\kappa_P=0$ and $|\varepsilon_P(c^{(0)})|\le Ct^{\theta_0/2}$ (Lemma~\ref{lem:NM}(2)). Since $\mathscr C_0$ is compact
and $c_0>0$ on $\mathcal Z$, there are $0<c_-\le c_+$ with $c^{(0)}_0\in[c_-,c_+]$ for all $c^{(0)}\in\mathscr C_0$. By
Lemma~\ref{lem:period}, with $\mu_P(c^{(0)})=1$, $\bigl||\mathcal B_q(c^{(0)})|-4\pi^2/c^{(0)}_0\bigr|\le Ct^{\theta_0/2}$, and
$|\mathcal B_q|\le b_+$ on $U$ with $b_+$ independent of $t$. So for $t$ small
\[
b_-:=2\pi^2/c_+\ \le\ |\mathcal B_q(c^{(0)})|\ \le\ b_+\qquad\text{for every }q .
\]
The polydisc in $\xi$ about $0$ of polyradius $\varrho_*=\min(1/2,\eta_0/8)$ lies in $U_\C\subseteq U$: on it
$\|\xi\|\le1/2$ and $|\eta_P|\le\|r_P\|_*\|\xi\|\le\eta_0/2$. Since $\mathcal B_q$ is holomorphic and bounded by $b_+$ on $U$,
the Cauchy estimates on this polydisc give $|\partial\mathcal B_q/\partial\xi_n(c^{(0)})|\le b_+/\varrho_*$. Hence, at $c^{(0)}$,
\[
d\wp_q=\mathcal B_q(c^{(0)})\,d\varepsilon_P+\varepsilon_P(c^{(0)})\,d\mathcal B_q
=\mathcal B_q(c^{(0)})\,r_P\cdot d\xi+O\bigl(C(\eta_0)\,t^{\theta_0/2}\bigr),
\]
by~\eqref{eq:depsilon} with $\eta_P=0$. Let $\mathbf D(t)$ be the $|Q|\times|\mathcal A|$ matrix whose rows are the
differentials $d\wp_q(c^{(0)})$, as linear forms in $d\xi$. Then
\begin{equation}\label{eq:Dmatrix}
\mathbf D(t)=\operatorname{diag}\bigl(\mathcal B_q(c^{(0)})\bigr)\,\mathbf R+\mathbf E(t),\qquad
\max_{q,n}|\mathbf E(t)_{qn}|\le C(\eta_0)\,t^{\theta_0/2}.
\end{equation}

\emph{Step 3: the rank.} Choose $r_{\mathbf R}$ rows $I$ and $r_{\mathbf R}$ columns $J$ of $\mathbf R$ with $\det\mathbf R_{IJ}\neq0$; as an integer,
$|\det\mathbf R_{IJ}|\ge1$. Then $\bigl|\det\bigl(\operatorname{diag}(\mathcal B)\mathbf R\bigr)_{IJ}\bigr|=\prod_{q\in I}|\mathcal B_q(c^{(0)})|\,
|\det\mathbf R_{IJ}|\ge b_-^{r_{\mathbf R}}$, and every entry of $\operatorname{diag}(\mathcal B)\mathbf R$ has modulus at most $b_+$. The determinant
is a polynomial of degree $r_{\mathbf R}$ in the entries, so, once $C(\eta_0)t^{\theta_0/2}\le1$,
\[
\bigl|\det\mathbf D(t)_{IJ}-\det\bigl(\operatorname{diag}(\mathcal B)\mathbf R\bigr)_{IJ}\bigr|
\le r_{\mathbf R}\cdot r_{\mathbf R}!\,(b_++1)^{r_{\mathbf R}-1}\,C(\eta_0)\,t^{\theta_0/2}.
\]
We decrease $t_0$ so that the right-hand side is less than $b_-^{r_{\mathbf R}}/2$. Then $\det\mathbf D(t)_{IJ}\neq0$, so
$\operatorname{rank}\mathbf D(t)\ge r_{\mathbf R}$, and the left kernel of $\mathbf D(t)$ in $\C^Q$ has dimension at most
$|Q|-r_{\mathbf R}=\operatorname{rank}K$.

\emph{Step 4: conclusion.} Let $k\in K_{\rm tor}$. Differentiating~\eqref{eq:periodrel} at $c^{(0)}$ gives
$\sum_qk_q\,d\wp_q(c^{(0)})=0$, so $k$ lies in the left kernel of $\mathbf D(t)$. The complex span of $K_{\rm tor}$ therefore has
dimension at most $\operatorname{rank}K$, and it contains the complex span of $K\subseteq K_{\rm tor}$, which has dimension
$\operatorname{rank}K$. The two spans agree, and since both $K_{\rm tor}$ and $K$ consist of integral vectors, $K_{\rm tor}\otimes\Q=K\otimes\Q$.
As $K$ is saturated in $\Z^Q$, $K_{\rm tor}\subseteq(K\otimes\Q)\cap\Z^Q=K$. Hence $K_{\rm tor}=\Rel(\delta)=K$.

The homomorphism $\Z^Q\to H_3(Y_t;\Z)$, $k\mapsto\sum_qk_q[\delta_q]$, therefore has kernel $K$, and its image meets the
torsion subgroup only in $0$. The image is isomorphic to $\Z^Q/K$, which is finitely generated and torsion-free because
$K$ is saturated, hence free of rank $|Q|-\operatorname{rank}K$. This proves item~(1). \qed

\begin{remark}\label{rem:relationsU}
Step~1 uses only that $\sum_qk_q[\delta_q]$ is torsion at one member, and the inclusion $K\subseteq\Rel(\delta)$ is the
integral statement of item~(2), not a statement modulo torsion. The rank bound of Step~3 is where the relation lattice
of the circuits enters: the left kernel of $\mathbf R$ over $\Q$ is the space of relations among the classes $[C_q]$ in
$H_2(\widehat X;\Q)$ (Proposition~\ref{prop:partialres}).
\end{remark}

The conclusion extends from the given member to all smooth members of $U$ that have Morse points at every node, as
Section~\ref{ssec:scope} states.

\begin{proposition}\label{prop:relationsU}
For every $c'\in U^\circ=U\setminus\bigcup_qD_q$, including complex coefficient vectors, the member $Y_{t,c'}$ is smooth, and,
with $\delta_q(c')$ oriented by transport from $c$, for $k\in\Z^Q$ the class $\sum_qk_q[\delta_q(c')]$ is a torsion element of $H_3(Y_{t,c'};\Z)$ if and only if $k\in K$, and
it is then zero.
\end{proposition}

\begin{proof}
Smoothness is Lemma~\ref{lem:U}(4). Parallel transport along a path in the connected set $U^\circ$ from $c$ to $c'$ is a
group isomorphism carrying $[\delta_q(c)]$ to $[\delta_q(c')]$ for every $q$ (Lemma~\ref{lem:flat}). It carries the
homomorphism $k\mapsto\sum_qk_q[\delta_q(c)]$ to $k\mapsto\sum_qk_q[\delta_q(c')]$ and the torsion subgroup onto the torsion
subgroup; so the kernels, and the preimages of the torsion subgroups, agree. Item~(1) at $c$ gives the claim.
\end{proof}

\subsection{The nodal locus: proof of Theorem \texorpdfstring{\ref{thm:main}}{A}(3)}\label{ssec:nodal}

By Lemma~\ref{lem:U}(4) the members of $U$ with a node at every point $q_P$, $P\in Q$, are the points of $D_Q$; each of
them has exactly $|Q|$ singular points, all ordinary double points. Let $I\subseteq Q$ index rows $r_P$, $P\in I$, that
form a basis of the row space of $\mathbf R$, so $|I|=r_{\mathbf R}=|Q|-\operatorname{rank}K$, and put $D_I:=\bigcap_{P\in I}D_P$. We
show that $D_I$ is a nonempty complex submanifold of $U$ of codimension $|I|$, and that $D_Q=D_I$.

In the chart $\xi$ write $\varepsilon_P(c^{(0)}e^\xi)=-\kappa_P(\xi)+E_P(\xi)$ with $\kappa_P=e^{-r_P\cdot\xi}-1$. By
Lemma~\ref{lem:NM}(2), $|E_P|\le Ct^{\theta_0/2}$ on $U^+_\C$. The space $\C^{\mathcal A}$, its subspaces and $\C^I$ carry
the max-norm, and operator norms are taken between these.

\emph{(a) $D_I\neq\emptyset$.} Let $V:=\operatorname{span}_\C\{r_P:P\in I\}\subseteq\C^{\mathcal A}$ and
${\mathsf A}\colon V\to\C^I$, ${\mathsf A}\xi:=(r_P\cdot\xi)_{P\in I}$. In the basis $(r_P)_{P\in I}$ of $V$ the matrix of ${\mathsf A}$ is the Gram
matrix of linearly independent real vectors, which is invertible; so ${\mathsf A}$ is an isomorphism, and $\|{\mathsf A}^{-1}\|\le C_{\mathsf A}$ for
a constant $C_{\mathsf A}$ depending only on $\mathbf R$ and $I$. Put $\varrho_V:=\min(1/2,\eta_0/8)$ and
$B_V:=\{\xi\in V:\|\xi\|\le\varrho_V\}$. Every point of $B_V$ lies in $U_\C\subseteq U$, and so does the polydisc of
polyradius $\varrho_*=\min(1/2,\eta_0/8)$ about it: on that polydisc $\|\xi\|\le1$ and $|\eta_P|\le\eta_0$. The Cauchy
estimates on these polydiscs give $\sup_{B_V}\|E\|\le Ct^{\theta_0/2}$ and $\sup_{B_V}\|dE\|\le C(\eta_0)t^{\theta_0/2}$
for $E:=(E_P)_{P\in I}$, with $C\le C(\eta_0)$. Let $t$ be so small that $C(\eta_0)t^{\theta_0/2}\le1/2$. On $B_V$ the
equations $\varepsilon_P=0$, $P\in I$, say $e^{-\eta_P}=1+E_P$. Here $|\eta_P|\le\eta_0/2<1$ and $|\log(1+E_P)|\le2|E_P|\le1$
for the principal logarithm, so $\eta_P=-\log(1+E_P)$, no other branch being possible, and the equations are equivalent
to the fixed point equation
\[
\xi=\mathsf T(\xi):={\mathsf A}^{-1}\bigl(-\log(1+E(\xi))\bigr),\qquad\xi\in B_V .
\]
For $|w|\le1/2$ one has $|\log(1+w)|\le2|w|$ and $|d\log(1+w)/dw|\le2$, and $B_V$ is convex. Hence
$\|\mathsf T(\xi)\|\le\varrho_t:=2C_{\mathsf A}\,C(\eta_0)\,t^{\theta_0/2}$, and $\mathsf T$ is Lipschitz on $B_V$ with constant at most
$2C_{\mathsf A}\sup_{B_V}\|dE\|\le\varrho_t$. We decrease $t_0$ so that $\varrho_t\le\varrho_V$ and $\varrho_t<1$. Then $\mathsf T$
maps the closed ball of radius $\varrho_t$ in $V$, which lies in $B_V$, into itself and is a contraction there. Its
fixed point $\xi^*$ satisfies $\varepsilon_P(c^{(0)}e^{\xi^*})=0$ for $P\in I$, with $\|\xi^*\|\le\varrho_t$, and
$c^{(0)}e^{\xi^*}\in U$. So $D_I\neq\emptyset$.

\emph{(b) $D_I$ is a complex submanifold of codimension $|I|$.} On $U$, $|e^{-\eta_P}|\in[e^{-1},e]$, and the
differentials~\eqref{eq:depsilon}, $P\in I$, are the rows of the matrix $\operatorname{diag}(e^{-\eta_P})_{P\in I}\mathbf R_I$ up to
an error bounded entrywise by $C(\eta_0)t^{\theta_0/2}$. The minor argument of Step~3 of Section~\ref{ssec:relations}, with
$e^{-\eta_P}$ in place of $\mathcal B_q$ and the bounds $e^{-1}$, $e$ in place of $b_\pm$, shows that these $|I|$ forms are
linearly independent at every point of $U$ once $t$ is small. By the holomorphic implicit function theorem $D_I$ is a
complex submanifold of $U$ of codimension $|I|$.

For use in (c) we bound $\eta_P$ on $D_I$. For $P\in I$, $\varepsilon_P=0$ gives $|e^{-\eta_P}-1|=|\kappa_P|=|E_P|\le
Ct^{\theta_0/2}$. On $U$, $|\eta_P|<\eta_0'=3\eta_0/2<1$, since $\eta_0\le r^2/20\le1/80$; and for $|w|\le1$ one has
$|e^{-w}-1|\ge|w|-(e-2)|w|^2\ge|w|/4$. Hence $|\eta_P|\le4Ct^{\theta_0/2}$ on $D_I$. For $P\notin I$,
$r_P=\sum_{P'\in I}a_{PP'}r_{P'}$ with rational coefficients $a_{PP'}$ depending only on $\mathbf R$ and $I$, so
$\eta_P=\sum_{P'}a_{PP'}\eta_{P'}$ and $|\eta_P|\le Ct^{\theta_0/2}$ on $D_I$ as well. Consequently
\begin{equation}\label{eq:epsDI}
|\varepsilon_P|\le|\kappa_P|+Ct^{\theta_0/2}\le Ct^{\theta_0/2}\qquad\text{on }D_I,\ \text{for every }P\in Q .
\end{equation}

\emph{(c) $D_Q=D_I$.} Let $q\notin I$, with parallelogram $P$. The rows $r_{P'}$, $P'\in I\cup\{P\}$, are linearly
dependent, and since the rows indexed by $I$ are independent they satisfy a rational relation $k$ with $k_q\neq0$ and
support in $I\cup\{q\}$. As $r_{P'}=-\circv_{q'}$, $k\in K\otimes\Q$, and a nonzero integer multiple $mk$ lies in $K$. By
item~(2) and~\eqref{eq:periodrel}, $\sum_{p}mk_p\wp_p\equiv0$ on $U$. On $D_I$, $\wp_p=\mathcal B_p\varepsilon_{P_p}=0$ for
$p\in I$, where $P_p$ is the parallelogram of $p$, so $k_q\mathcal B_q\varepsilon_P=0$ there. By~\eqref{eq:epsDI}, $|\varepsilon_P|\le\varepsilon_*$ on $D_I$ once $t$ is
small, so $\mathcal B_q$ has no zero on $D_I$ (Lemma~\ref{lem:period}). Hence $\varepsilon_P=0$ on $D_I$, that is
$D_I\subseteq D_q$. So $D_I\subseteq D_Q\subseteq D_I$, and $D_Q=D_I$ is a nonempty complex submanifold of $U$ of
codimension $|Q|-\operatorname{rank}K$.

Finally, by Proposition~\ref{prop:partialres}(3) the diagonals can be chosen so that the small resolution
$\widehat X\to X'_\sigma$ is projective, and for every projective small resolution $|Q|-\operatorname{rank}K$ is the relative
Picard number of $\widehat X\to X'_\sigma$ (Proposition~\ref{prop:partialres}(2); changing the diagonal at $q$ changes only the
sign of the image of $[C_q]$ in $\Q^R$, and by Lemma~\ref{lem:diag} it leaves items~(1) and~(2) unchanged). This proves item~(3), and with it Theorem~\ref{thm:main}. \qed

\subsection{Lagrangian representatives}\label{ssec:lagrangian}

Casta\~no-Bernard and Matessi construct Lagrangian torus fibrations over integral affine three-manifolds with
singularities \cite{CBM09}, and Lagrangian submanifolds associated with tropical varieties are constructed in
\cite{Mik19,Mat21a,Mat21b}; here the Lagrangian spheres are built directly in the K\"ahler manifolds $Y_t$. Fix a toric
K\"ahler form $\omega_{\mathcal T}$ on the smooth projective toric variety $X_{\mathcal T}$, and write $\omega_{c'}$ for its
restriction to a smooth member $Y_{t,c'}$, a K\"ahler form.

\begin{lemma}\label{lem:L}
Let $c$ be a member as in Theorem~\ref{thm:main}, with neighbourhood $U$, and let $Q'\subseteq Q$ be a set of nodes with
$D_{Q'}\neq\emptyset$. Then there are pairwise disjoint embedded Lagrangian three-spheres $L_P\subset(Y_t,\omega_c)$,
$P\in Q'$, each smoothly isotopic in $Y_t$ to $\delta_P$.
\end{lemma}

The lemma controls the isotopy class of $L_P$, without a bound on its position in the Morse polydisc: for a fixed K\"ahler form the equivalence constant between
$\omega_{\mathcal T}$ and the flat metric of the rescaled Morse chart grows like a negative power of $t$, and at the member
$c$ itself the thimbles need not stay in the polydiscs.

\begin{proof}
\emph{Step 1: a thimble bound near $D_{Q'}$.} Fix $b_0\in D_{Q'}$ and a compact neighbourhood $\mathcal V\subset U$ of
$b_0$. By Lemma~\ref{lem:NM}, which holds on $\mathcal V$ including where $\varepsilon_P=0$, each $f_P:=f^{(1)}$ has on
$\mathbb D_P$ a single critical point $q_P$, and it is nondegenerate. By the holomorphic Morse lemma with parameters
\cite[\S\S6.2, 9.6]{AGZV1} there are holomorphic coordinates $\zeta$ centred at $q_P$, depending continuously on the member, with
$f_P=\varepsilon_P+\sum_i\zeta_i^2$ exactly. Put $\Phi_P:=\operatorname{Re}(e^{-i\arg\varepsilon_P}f_P)$; then $|d\Phi_P|_0=2|\zeta|$,
where $|\cdot|_0$ is the flat norm of the coordinates $\zeta$. By compactness there are $\varpi_0>0$ and $C\ge1$, uniform on
$\mathcal V$, such that the ball $\mathsf B_P:=\{|\zeta|\le\varpi_0\}$ lies in the chart and in $\mathbb D_P$, and
$C^{-1}|\cdot|_0\le|\cdot|_\omega\le C|\cdot|_0$ on $\mathsf B_P$, where $|\cdot|_\omega$ is the norm of
$\omega_{\mathcal T}$. Along a flow line of $-\nabla_\omega\Phi_P$, $\Phi_P$ decreases at rate $|d\Phi_P|_\omega\ge C^{-1}|d\Phi_P|_0$ per unit
$\omega$-length, and the $\omega$-length needed to cross the shell $\varpi/2\le|\zeta|\le\varpi$ is at least $\varpi/(2C)$;
so crossing it costs a drop of $\Phi_P$ of at least $c_H\varpi^2/(4C^2)$ with $c_H=2$, while the drop from $q_P$ to the level
$f_P=0$ is $|\varepsilon_P|$. Hence, if $|\varepsilon_P|<\Phi_0:=c_H\varpi_0^2/(4C^2)$, the descending manifold of $q_P$ between
the levels $\varepsilon_P$ and $0$ lies in the ball of radius $2C(|\varepsilon_P|/c_H)^{1/2}$, and every flow line from $q_P$
reaches the level $0$ there, $q_P$ being the only critical point.

\emph{Step 2: Lagrangian spheres near $D_{Q'}$.} Choose $b_1\in\mathcal V\cap U^\circ$ with $0<|\varepsilon_P(b_1)|<\Phi_0$
for $P\in Q'$; this is possible because the $\varepsilon_P$ are continuous, vanish at $b_0$, and $\bigcup_qD_q$ is nowhere
dense. By the Cauchy--Riemann equations the gradient flow of $\Phi_P$ for the K\"ahler metric preserves
$\operatorname{Im}(e^{-i\arg\varepsilon_P}f_P)$, so by Step~1 the descending manifold is the Lefschetz thimble of $q_P$ for
the straight path from $\varepsilon_P$ to $0$, contained in $\mathsf B_P$. By \cite[\S1.3, Lemmas~1.13--1.14]{Seidel},
whose argument is local along the thimble, it is an embedded Lagrangian ball, and its boundary $L_P(b_1)$ is an embedded
Lagrangian sphere in $Y_{t,b_1}$. In our setting the Lagrangian property is seen directly: by the Cauchy--Riemann
equations the gradient of $\Phi_P$ is a Hamiltonian vector field of $\operatorname{Im}(e^{-i\arg\varepsilon_P}f_P)$, so its
flow preserves $\omega_{\mathcal T}$; it carries the points of the thimble to $q_P$ as time tends to infinity and contracts
the tangent spaces of the thimble; hence these tangent spaces are isotropic, and the thimble is Lagrangian. On the thimble
$\operatorname{Im}(e^{-i\arg\varepsilon_P}f_P)=0$, so its boundary is the regular level $\{\Phi_P=0\}$ of $\Phi_P$ on it, an
isotropic three-sphere in $Y_{t,b_1}$, hence Lagrangian; the flow lines reach this level by Step~1. Seidel credits the
definition of the thimble to a suggestion of Donaldson; compare \cite{Donaldson}. The
polydiscs are disjoint, so the $L_P(b_1)$ are. Replacing $\omega_{\mathcal T}$ on $\mathsf B_P$ by
$(1-\tau)\omega_{\mathcal T}+\tau\,\mathcal X_*\omega_{\rm std}$, $\tau\in[0,1]$, where $\mathcal X$ is the Morse chart
$\zeta$ of Step~1, keeps the estimates of Step~1 with the same constant $C$, since $C\ge1$, and gives a smooth isotopy in
$\mathbb D_P\cap Y_{t,b_1}$ from $L_P(b_1)$ to the Morse-chart sphere $\delta_P(b_1)$.

\emph{Step 3: symplectic parallel transport.} Choose a smooth path $\gamma$ in the connected set $U^\circ$ from $b_1$ to
$c$. Every $Y_{t,\gamma(s)}$ is smooth, so $\mathcal Y=\{(s,y):y\in Y_{t,\gamma(s)}\}$ is a manifold submersing onto
$[0,1]$ with compact fibres. Let $\omega_{\mathcal Y}$ be the pullback of $\omega_{\mathcal T}$. Since the fibres are
symplectic, their $\omega_{\mathcal Y}$-orthogonal complement is a line field transverse to them; let $\hat\partial_s$ be
the lift of $\partial_s$ in it, whose flow $\phi_s$ is defined on $[0,1]$. Then
$\mathcal L_{\hat\partial_s}\omega_{\mathcal Y}=d(\iota_{\hat\partial_s}\omega_{\mathcal Y})$, and
$\iota_{\hat\partial_s}\omega_{\mathcal Y}$ vanishes on the fibres, so each $\phi_s$ restricts to a symplectomorphism between
the fibres. Put $L_P:=\phi_1(L_P(b_1))$: these are pairwise disjoint embedded Lagrangian spheres in $(Y_t,\omega_c)$.

\emph{Step 4: the isotopy to $\delta_P$.} Along $\gamma$, $\varepsilon_P\ne0$ and Lemma~\ref{lem:NM} holds, so by the
holomorphic Morse lemma with parameters the Morse-chart spheres $\delta_P(\gamma(s))$ form a smooth family. Two holomorphic
Morse charts at $q_P$ differ by a biholomorphism $g$ preserving $\sum_i\zeta_i^2$; the maps $\zeta\mapsto g(\lambda\zeta)/\lambda$,
$\lambda\in(0,1]$, preserve it as well and tend to $dg(0)\in O(4,\C)$, and $SO(4,\C)$ is connected while $O(4,\R)$ preserves
the real sphere, so the sphere is well defined up to isotopy. Pulling the spheres $\delta_P(\gamma(s))$ back by $\phi_s$
gives an isotopy in $Y_{t,b_1}$ from $\delta_P(b_1)$ to $\phi_1^{-1}(\delta_P(c))$. Composing with the isotopy of Step~2
and applying $\phi_1$ gives a smooth isotopy in $Y_t$ from $L_P$ to $\delta_P$. We orient $L_P$ by transport along
this isotopy from $\delta_P=\delta_P(c)$, which carries the orientation of Definition~\ref{def:orient}; along $\gamma$ the
spheres $\delta_P(\gamma(s))$ carry the orientation by transport from $c$ of Section~\ref{ssec:U}.
\end{proof}

By Theorem~\ref{thm:main}(3), $D_Q\neq\emptyset$, so $D_{Q'}\supseteq D_Q$ is nonempty for every $Q'\subseteq Q$, and
Lemma~\ref{lem:L} applies with $Q'=Q$: the spheres $\delta_q$, $q\in Q$, are simultaneously isotopic to pairwise disjoint
Lagrangian spheres $L_q$, with $[L_q]=[\delta_q]$ in $H_3(Y_t;\Z)$.

\subsection{Smoothings: proof of Corollary \texorpdfstring{\ref{cor:smooth}}{B}}\label{ssec:smooth}

The equivalence of (a) and (b) is Namikawa's theorem \cite[Thm.~2.5]{Nam02}, which we apply after verifying its
hypotheses. Namikawa calls a projective threefold with only terminal singularities, trivial canonical class and
$H^1(\mathcal O)=0$ a \emph{Calabi--Yau threefold} \cite[\S0.2]{Nam02}.

\begin{lemma}\label{lem:Xprime}
$X'_\sigma$ is a Calabi--Yau threefold in the sense of \cite{Nam02}: it is a normal projective threefold whose singular
points are the $|Q|$ nodes, with $\omega_{X'_\sigma}\cong\mathcal O_{X'_\sigma}$ and
$H^1(X'_\sigma,\mathcal O)=H^2(X'_\sigma,\mathcal O)=0$. The small resolution $\widehat X\to X'_\sigma$ is an isomorphism off the
curves $C_q\cong\mathbb P^1$.
\end{lemma}

\begin{proof}
The toric variety $P'_\sigma$ of~\ref{hyp:P} is projective, and it is Gorenstein: its rays are generated by lattice points
of $\partial\Delta$, so the support function $\varphi$ of $-K$ is linear on each of its cones. The morphism
$P'_\sigma\to\mathbb P_\Delta$ is crepant for the same reason. So the total transform of the anticanonical hypersurface $X$
is an anticanonical divisor of $P'_\sigma$. It has no exceptional component: such a component would map onto the closure
of a torus orbit of $\mathbb P_\Delta$ contained in $X$, while the $\Delta^\circ$-regular hypersurface $X$ meets every orbit
transversally and contains none \cite[\S3]{Batyrev}. Hence $X'_\sigma\in|-K_{P'_\sigma}|$ is a Cartier divisor, and it is
projective. By adjunction $\omega_{X'_\sigma}\cong\mathcal O_{X'_\sigma}$. The singular points of $X'_\sigma$ are the nodes,
one for each $q\in Q$ (Section~\ref{ssec:partialres}); ordinary double points are terminal. As a Cartier divisor in the
Cohen--Macaulay variety $P'_\sigma$, $X'_\sigma$ is Cohen--Macaulay, and with finitely many singular points it is normal.
From $0\to\omega_{P'_\sigma}\to\mathcal O_{P'_\sigma}\to\mathcal O_{X'_\sigma}\to0$, Serre duality on the four-dimensional
Cohen--Macaulay projective variety $P'_\sigma$, and $H^j(P'_\sigma,\mathcal O)=0$ for $j>0$ \cite[\S9.2]{CLS}, we get
$H^i(P'_\sigma,\omega_{P'_\sigma})\cong H^{4-i}(P'_\sigma,\mathcal O)^\vee=0$ for $i\le3$, hence
$H^0(X'_\sigma,\mathcal O)=\C$ and $H^i(X'_\sigma,\mathcal O)=0$ for $i=1,2$; in particular $X'_\sigma$ is connected, hence
irreducible. The last sentence is the description of $\widehat X$ in Section~\ref{ssec:partialres}.
\end{proof}

\begin{lemma}\label{lem:allnonzero}
Let $L\subseteq\Z^Q$ be a subgroup. The following are equivalent: (i) some $\lambda\in L$ has $\lambda_q\neq0$ for every
$q$; (ii) some $\lambda\in L\otimes\Q$ does; (iii) some $\lambda\in L\otimes\C$ does; (iv) for no $q$ is $L$ contained in the
coordinate hyperplane $\{\lambda_q=0\}$.
\end{lemma}

\begin{proof}
(i)$\Rightarrow$(ii)$\Rightarrow$(iii) holds because $L\subseteq L\otimes\Q\subseteq L\otimes\C$, and (iii)$\Rightarrow$(iv) because $L\subseteq\{\lambda_q=0\}$ implies
$L\otimes\C\subseteq\{\lambda_q=0\}$. For (iv)$\Rightarrow$(i) choose $\ell^{(p)}\in L$ with $\ell^{(p)}_p\neq0$ for each
$p\in Q$. For $m\in\Z^Q$, the $q$-th coordinate of $\sum_pm_p\ell^{(p)}$ is a linear form in $m$ whose coefficient of
$m_q$ is $\ell^{(q)}_q\neq0$. The union of the $|Q|$ kernels of these nonzero forms does not contain $\Z^Q$, so some
$\lambda=\sum_pm_p\ell^{(p)}$ has every coordinate nonzero.
\end{proof}

\begin{proof}[Proof of Corollary~\ref{cor:smooth}]
\emph{(a)$\Leftrightarrow$(b).} By Lemma~\ref{lem:Xprime}, $X'_\sigma$ is a Calabi--Yau threefold in Namikawa's sense
whose singular points are the nodes $q\in Q$, and $\widehat X\to X'_\sigma$ is a small resolution, not necessarily projective,
that is an isomorphism off the curves $C_q\cong\mathbb P^1$. For these data \cite[Thm.~2.5]{Nam02} states that $X'_\sigma$ is
smoothable by a flat deformation if and only if some relation $\sum_q\lambda_q[C_q]=0$ in $H_2(\widehat X;\C)$ has every
$\lambda_q\neq0$. (The theorem adds a third equivalent condition, maximality of $X'_\sigma$, which we do not use.
Friedman's criterion \cite[Rem.~(4.5), Cor.~(4.7)]{Friedman} gives the sufficiency of such a relation under the
hypotheses of \cite[Thm.~(4.4)]{Friedman}, which include that the exceptional curve classes span
$H^2(\widehat X,\Omega^2_{\widehat X})$;
Namikawa's theorem does not require that hypothesis, and gives both directions for flat smoothings.) The classes $[C_q]$ are rational, and by
Proposition~\ref{prop:partialres} the relations among them in $H_2(\widehat X;\Q)$ form $K\otimes\Q$; the relations in
$H_2(\widehat X;\C)=H_2(\widehat X;\Q)\otimes\C$ therefore form $K\otimes\C$. By Lemma~\ref{lem:allnonzero} applied to $L=K$,
a relation with every coefficient nonzero exists over $\C$ if and only if one exists over $\Q$. This is (b).

\emph{(b)$\Leftrightarrow$(c).} By Proposition~\ref{prop:partialres} and Lemma~\ref{lem:allnonzero}, (b) holds if and only if
some $\lambda\in K$ has every $\lambda_q\neq0$. By Theorem~\ref{thm:main}(1), $K=\Rel(\delta)$, the lattice of relations
among the classes $[\delta_q]$ in $H_3(Y_t;\Z)$; so this holds if and only if (c) holds. Conditions (b) and (c) do not
depend on the diagonals: a change of diagonal at $q$ changes the sign of $[C_q]$ (Proposition~\ref{prop:partialres}(2)) and
of $\delta_q$ (Lemma~\ref{lem:diag}), which preserves the existence of a relation with every coefficient nonzero. By
Proposition~\ref{prop:relationsU} the same equivalence holds at every member of $U^\circ$.

\emph{The symplectic statement.} By Theorem~\ref{thm:main}(3) and Lemma~\ref{lem:L} with $Q'=Q$
(Section~\ref{ssec:lagrangian}), the spheres $\delta_q$ are smoothly isotopic to pairwise disjoint embedded Lagrangian
spheres $L_q\subset(Y_t,\omega_c)$, $q\in Q$, with $[L_q]=[\delta_q]$; this part is unconditional and does not use
(a)--(c). The manifold $Y_t$
is a compact K\"ahler, in particular symplectic, six-manifold. If (c) holds, $\sum_qk_q[L_q]=0$ in $H_3(Y_t;\Z)$ with every
$k_q\neq0$, which is the hypothesis of \cite[Thm.~2.9]{STY} for the $n=|Q|$ spheres $L_q$: a relation in $H_3(\,\cdot\,;\Z)$
among $n$ disjoint embedded Lagrangian three-spheres in a symplectic six-manifold, with every coefficient nonzero. That
theorem gives a symplectic structure on one of the $2^{|Q|}$ conifold transitions of $Y_t$ in the spheres $L_q$ in which
the exceptional curves are symplectic. In the proof of \cite[Thm.~2.9]{STY} the choice at the sphere $L_q$, an orientation
of $L_q$ together with one of the two small resolutions, is determined by the sign of $k_q$.
\end{proof}

\begin{remark}[forced nodes]\label{rem:forced}
By Lemma~\ref{lem:allnonzero}, (a)--(c) fail exactly when some node $q$ is \emph{forced}: every $\lambda\in K$ has
$\lambda_q=0$. By Theorem~\ref{thm:main}(1) this is the case if and only if $\delta_q$ occurs with nonzero coefficient in
no relation among the classes $[\delta_{q'}]$, that is, if and only if $[\delta_q]$ does not lie in the span of the other
classes in $H_3(Y_t;\Q)$. The proof of \cite[Thm.~2.5]{Nam02} shows that a forced
node of $X'_\sigma$ survives every flat deformation. Whether (b) holds is a finite linear-algebra test on the circuits
$\circv_q$; Section~\ref{ssec:list} records its outcome on the list.
\end{remark}

\begin{corollary}[smoothings of $X$ with $dP_7$ points]\label{cor:mixed}
Let $\Delta$ be admissible with no hexagonal two-face, so that the profile $\sigma$ is unique, and assume~\ref{hyp:P}. Let
$Y_t$ be a member as in Theorem~\ref{thm:main}. Then the following are equivalent:
\begin{enumerate}
\item[(a)] $X$ admits a smoothing;
\item[(b)] $X'_\sigma$ admits a smoothing;
\item[(c)] the vanishing spheres satisfy a relation $\sum_{q\in Q}k_q[\delta_q]=0$ in $H_3(Y_t;\Z)$ with $k_q\neq0$ for
every $q\in Q$, including both nodes of every pentagon.
\end{enumerate}
In~(c), with the diagonals through the centres chosen after Definition~\ref{def:profile}, $k_{q_1}+k_{q_2}=0$ at the two
nodes $q_1,q_2$ of each pentagon.
\end{corollary}

\begin{proof}
\emph{$X$ satisfies the hypotheses of \cite[Thm.~4.7]{Paper5}}, which are stated in \cite[\S1]{Paper5}. The argument of
Lemma~\ref{lem:Xprime}, with $\mathbb P_\Delta$ in place of $P'_\sigma$, shows that $X$ is a connected normal projective
threefold with $\omega_X\cong\mathcal O_X$ and $H^1(X,\mathcal O_X)=0$. By Section~\ref{ssec:partialres} its singular points
are the nodes $x_G$ of the squares and, at each pentagon $G$, a point analytically isomorphic to the vertex of $U_G$, the
anticanonical cone over the toric surface of $G$, which is $dP_7$. So \cite[Thm.~4.7]{Paper5} applies. At a $dP_7$ point
the local base is $\C\{s,w\}/(w^2,sw)$ \cite[\S2.1]{Paper5}, whose reduction is a line; so there is exactly one branch
profile $S$ \cite[Def.~2.1]{Paper5}, and every one-parameter deformation over a disc induces it, since a morphism from a
reduced space factors through the reduction of the target. A smoothing in our sense, a flat deformation over a germ of
complex space whose general fibre is smooth, gives a convergent one-parameter smoothing, that is, a flat proper family over
a disc with smooth fibres off the centre: by \cite[Cor.~6.7]{Paper4}, whose hypotheses (a compact threefold with trivial
canonical sheaf, $H^1(\mathcal O)=0$, only nodes and anticanonical cones over $dP_6$ or $dP_7$ as singularities, and a crepant resolution, small over the nodes
and with an exceptional del Pezzo surface over each cone point) $X$ satisfies by the above and by the resolution
$\widehat X\to X$ constructed below, a threefold of this kind without a one-parameter smoothing is not smoothable over any base germ;
its proof applies the curve selection lemma to the discriminant in the base of a miniversal deformation. The converse is
trivial. So (a) holds if and only if the criterion of
\cite[Thm.~4.7]{Paper5} holds: the space $K_S=\ker\bigl(\bigoplus_pR_{p,S_p}\to H_2(Y_0;\Q)\bigr)$ of
\cite[\S1, Def.~2.1]{Paper5}, where every summand $R_{p,S_p}$ is a line, contains a vector with nonzero component at every
node and at every $dP_7$ point.

\emph{The resolution.} Let $\widehat X$ be the small resolution of $X'_\sigma$ with the chosen diagonals, which pass
through the centres at the pentagons. Near the fibre over a pentagon point it is the toric variety of the star subdivision
of the cone over $G$, that is, a neighbourhood of the zero section of $\operatorname{Tot}(K_E)$ with
$E=D_o\cong\mathrm{Bl}_2\mathbb P^2$; at the squares it is a small resolution of the node. So $\widehat X\to X$ is the
resolution $Y_0\to X$ of \cite[\S2.1]{Paper5}. By Proposition~\ref{prop:classical}(1), $C_{q_1}=e_1$ and $C_{q_2}=e_2$ are
two disjoint $(-1)$-curves of $E$. The line $R_{p,S_p}=\Q(e_1-e_2)\subset K_E^\perp=H_2(L_p;\Q)$ of \cite[\S2.1]{Paper5}
is spanned by the $(-2)$-classes of $K_E^\perp$, which are $\pm(e_1-e_2)$ in every marking, so it does not depend on the
marking, and its image in $H_2(\widehat X;\Q)$ is spanned by $[C_{q_1}]-[C_{q_2}]$. At a node, $R_{p,S_p}=\Q[C_q]$.

\emph{The dictionary.} Define $\Phi\colon K_S\to\Q^Q$ by $\Phi(\kappa)_q=\kappa_q$ at a square and
$\Phi(\kappa)_{q_1}=\kappa_p$, $\Phi(\kappa)_{q_2}=-\kappa_p$ at a pentagon $p$, where $\kappa_p(e_1-e_2)$ is the component
of $\kappa$ at $p$. Then $\sum_q\Phi(\kappa)_q[C_q]$ is the image of $\kappa$ in $H_2(\widehat X;\Q)$, which vanishes. By
Proposition~\ref{prop:partialres}(2), $\Phi(K_S)\subseteq K\otimes\Q$. Conversely, every $\lambda\in K\otimes\Q$ satisfies
$\lambda_{q_1}+\lambda_{q_2}=0$ at each pentagon (Remark~\ref{rem:tie}), so it equals $\Phi(\kappa)$ for the $\kappa$ with
$\kappa_q=\lambda_q$ at the squares and $\kappa_p=\lambda_{q_1}$ at the pentagons. So $\Phi$ is an isomorphism onto
$K\otimes\Q$. The components of $\kappa$ are all nonzero if and only if every coordinate of $\Phi(\kappa)$ is nonzero, both
nodes of each pentagon included. By Lemma~\ref{lem:allnonzero}, (a) is therefore equivalent to the existence of
$\lambda\in K$ with every $\lambda_q\neq0$.

\emph{Conclusion.} By Corollary~\ref{cor:smooth} and Lemma~\ref{lem:allnonzero}, the same condition is equivalent to (b).
By Theorem~\ref{thm:main}(1), $K=\Rel(\delta)$, so it is equivalent to (c). The tie holds on $\Rel(\delta)=K$.
\end{proof}

\subsection{The factoring locus and the nodal locus}\label{ssec:periodsremarks}

Batyrev and Kreuzer specialised the mirror family, for polytopes whose singular two-faces are squares, to the locus where
the face polynomial of each square is a product of two binomials \cite[\S3]{BK10}; we call its analogue for all node
parallelograms, $\mu_P=1$ for every $P$, the \emph{factoring locus}. They expected the specialised members to
have one conifold singularity at the point $q_j$ of each of their local affine toric charts $U_{\sigma_j}$, and
derived it from a local equation in which every monomial other than those of the square is divisible by the product
of the two normal coordinates; they remark that this argument does not show that these are all the singular points
\cite[\S3]{BK10}. In the node chart of Section~\ref{ssec:morse}, where the normal coordinates are $Z$ and $Z'$, that local
equation omits the terms $AZ+A'Z'$ of the apexes, which are linear in one normal coordinate (for dual edges of lattice
length one, the case considered here). These terms are present: the
three-cells $\omega_1,\omega_2$ containing $P$ have apexes, and by Proposition~\ref{prop:pulling}(4) every apex has
exponent one in exactly one of $Z$, $Z'$ and exponent zero in the other. Within $U$ the factoring locus
is $\{c'\in U:\eta_P(c')=0\text{ for every }P\}$, because $\mu_P=e^{\eta_P}$ and $|\eta_P|<1$ on $U$; its points with real
coefficients of our sign pattern form $\mathcal Z\cap U$. The nodal locus $D_Q$ lies close to it: by the bound on the
$\eta_P$ in the proof of item~(3)(b), $|\eta_P|\le Ct^{\theta_0/2}$ on $D_Q$ for every $P$.

The estimates do not place $D_Q$ on the factoring locus. On the factoring locus $\kappa_P=0$, so $\varepsilon_P=E_P$ in the
notation of Section~\ref{ssec:nodal}; Lemma~\ref{lem:NM}(2) bounds $E_P$ by $Ct^{\theta_0/2}$, but the estimates do not make it
zero, and by Lemma~\ref{lem:U}(4) the member is singular near $q_P$ only if $\varepsilon_P=0$. Among the terms that move the
critical value are those of the apexes of the cells next to $P$; at a pentagon or hexagon node the monomials of the other
lattice points of the two-face, which are constant in $Z$ and $Z'$, move it as well. At a square node, in the node chart
at $\kappa_P=0$, the monomials of $\mathcal A_P$ give $xy+\mathfrak b\,ZZ'+AZ+A'Z'$, with the apex sums $A$ and $A'$. In
the local model where these are replaced by $A_0+A_1x$ and $A'_0+A'_1y$, with constants $A_0,A_1,A'_0,A'_1$ and
$A_1A'_1\neq\mathfrak b$, the critical point is $Z=-A'_0/(\mathfrak b-A_1A'_1)$, $Z'=-A_0/(\mathfrak b-A_1A'_1)$,
$x=-A'_1Z'$, $y=-A_1Z$, and the critical value is
\[
-\frac{A_0A'_0}{\mathfrak b-A_1A'_1},
\]
which is nonzero whenever the apex sums do not vanish at the node point. So the estimates do not show that the members of
the factoring locus are nodal at $q_P$, and in this local model at a square node they are not, unless an apex sum vanishes at the node point; the members of $U$ with a
node at every $q_P$ are those of $D_Q$. For these,
Corollary~\ref{cor:smooth} settles the analogue of the question left open in \cite[\S3]{BK10}: the spheres $\delta_q$
satisfy a relation with every coefficient nonzero exactly when $X'_\sigma$ is smoothable. The corollary produces a
symplectic conifold transition, by \cite[Thm.~2.9]{STY}, and not a projective small resolution.

Theorem~\ref{thm:main}(3) is the mirror half of Morrison's count (Section~\ref{sec:statements}) at the members of $U$. The
other half of the proposal in \cite[\S3.1]{Mor99}, that the members of $D_Q$ have a small resolution mirror to a smoothing
of $X'_\sigma$, is not addressed here.
